% ════════════════════════════════════════════════════════════════
%  INTRODUCTION GÉNÉRALE
% ════════════════════════════════════════════════════════════════
\chapter*{Introduction générale}

\section*{Contexte du projet}
\addcontentsline{toc}{section}{Contexte du projet}

La fonction Ressources Humaines constitue un levier direct de performance
dans un environnement industriel exigeant. Chez
\textbf{Schulte Automotive Tunisia Sarl}, le recrutement concerne deux
sites de production (Bouarada et Zaghouan), avec des besoins continus en
profils techniques et opérateurs.

Avant ce projet, le traitement des candidatures reposait principalement sur
des pratiques manuelles (papier, courrier, échanges non centralisés). Ce
mode de fonctionnement générait des délais, un risque de perte
d'information et une visibilité limitée sur l'avancement des dossiers,
autant pour les candidats que pour l'équipe RH.

\section*{Problématique}
\addcontentsline{toc}{section}{Problématique}

Le problème central est l'absence d'une plateforme locale de recrutement
capable de publier les offres, de centraliser les candidatures et de
suivre chaque étape du processus de manière traçable.

En pratique, trois difficultés ressortent :
\begin{itemize}
    \item un accès peu fluide aux offres pour les candidats ;
    \item un suivi interne difficile pour le service RH ;
    \item une communication insuffisante sur l'état d'avancement des candidatures.
\end{itemize}

\section*{Objectifs}
\addcontentsline{toc}{section}{Objectifs}

Ce projet vise à concevoir et réaliser une solution numérique de
recrutement, adaptée au contexte de Schulte Tunisia.

\textbf{Objectif général :} mettre en place un processus de recrutement
numérique, structuré et transparent.

\textbf{Objectifs spécifiques :}
\begin{itemize}
    \item digitaliser la publication et la gestion des offres d'emploi ;
    \item centraliser les candidatures dans un espace unique ;
    \item améliorer la traçabilité des décisions RH ;
    \item offrir au candidat un suivi clair de sa candidature ;
    \item réduire les tâches administratives répétitives.
\end{itemize}

Au-delà de la simple digitalisation, la solution vise aussi à garantir une
cohérence technique entre plusieurs interfaces clientes, une sécurité
adaptée aux rôles, et une traçabilité fiable des actions métier.

\section*{Méthodologie adoptée (aperçu rapide)}
\addcontentsline{toc}{section}{Méthodologie adoptée (aperçu rapide)}

La démarche suivie combine analyse métier et développement itératif.
Nous avons d'abord étudié le processus existant et formalisé les besoins
fonctionnels et non fonctionnels. La réalisation a ensuite été structurée
en sprints Scrum, avec des incréments successifs et des validations
régulières avec les parties prenantes.

Cette méthode a permis d'avancer par cycles courts, de corriger rapidement
les écarts constatés et de maintenir un alignement constant avec les
attentes terrain.

\section*{Structure du rapport}
\addcontentsline{toc}{section}{Structure du rapport}

Le rapport est organisé en cinq chapitres complémentaires :
\begin{enumerate}
    \item Le premier chapitre présente l'organisme d'accueil, le cadre du
    projet, la problématique et la solution proposée.

    \item Le deuxième chapitre expose l'état de l'art, l'étude comparative
    des solutions existantes et les choix technologiques retenus.

    \item Le troisième chapitre présente l'analyse et la spécification
    des besoins, les acteurs du système et la modélisation UML.

    \item Le quatrième chapitre présente la conception et la réalisation
    de la solution, organisée en releases Scrum.

    \item Le cinquième chapitre expose la stratégie de test, les cas
    exécutés, les résultats obtenus et la validation du système.
\end{enumerate}

Le rapport se termine par une conclusion générale et des perspectives
d'amélioration.
